\documentclass{article}
\usepackage[utf8]{inputenc}
\usepackage[hidelinks]{hyperref}
\usepackage{abstract}
\usepackage{lipsum}
\usepackage{hyperref}
\usepackage{verbatim}
\usepackage{graphicx}
\usepackage{float}
\usepackage{indentfirst}
\usepackage{listings}
\usepackage[acronym]{glossaries}
\usepackage{cleveref}

\bibliographystyle{ieeetr}

\graphicspath{ {./images/} }

\crefname{figure}{\textbf{Figure.}}{\textbf{Figures.}}
\crefname{table}{\textbf{Table.}}{\textbf{Tables.}}

\pagenumbering{roman}

\makeglossaries
\newacronym{iot}{IoT}{Internet of Things}
\newacronym{m2m}{M2M}{Machine-to-Machine}

\renewcommand{\contentsname}{\centering TABLE OF CONTENTS}
\renewcommand{\listtablename}{\centering LIST OF TABLES}
\renewcommand{\listfigurename}{\centering LIST OF FIGURES}

\renewcommand{\abstractname}{}
\renewcommand{\absnamepos}{empty}

\providecommand{\keywords}[1]
{
  \small
  \textbf{\textit{Key Words---}} #1
}

\title{
	\Large
    T.C.\\ BAHÇEŞEHİR UNIVERSITY
    \includegraphics[width=\textwidth]{images/bau-logo.jpeg}
    
    \large 
    FACULTY OF ENGINEERING AND NATURAL SCIENCES\\[1cm]
    CAPSTONE FINAL REPORT\\[1cm]
    THE TITLE OF THE PROJECT GOES HERE\\[1cm]
}
\author {
    The Names of the Students and their Departments go here,\\
    In Alphabetical Order by Surname,\\
    Capitalize The First Letters Only\\[1cm]
    Advisors: Your Advisor’s Name Goes Here\\
    (e.g. Assist. Prof., Assoc. Prof., Prof. ...) and their departments\\[1cm]
}
\date{ISTANBUL, May 2022}

\begin{document}
\maketitle
\thispagestyle{empty}

\newpage
\section*{\centering STUDENT DECLARATION}
\setcounter{page}{2}
By submitting this report, as partial fulfillment of the requirements of the Capstone course, the students promise on penalty of failure of the course that
\begin{itemize}
    \item they have given credit to and declared (by citation), any work that is not their own (e.g. parts of the report that is copied/pasted from the Internet, design or construction performed by another person, etc.);
    \item they have not received unpermitted aid for the project design, construction, report or presentation;
    \item they have not falsely assigned credit for work to another student in the group, and not take credit for work done by another student in the group. 
\end{itemize}
\addcontentsline{toc}{section}{STUDENT DECLARATION}

\newpage
\section*{\centering ABSTRACT}
\begin{center}
    THE TITLE OF THE PROJECT GOES HERE\\[0.25cm]
    Your Names in alphabetical order by surname\\[0.75cm]
    Faculty of Engineering and Natural Sciences\\[0.25cm]
    Advisor: Your Adviser Titles and Names\\[0.25cm]
    May 2021\\[1cm]
\end{center}

\begin{abstract}
    Write the abstract here (only necessary for the final report).\\\
    \par
    \lipsum[1]
    \par
    \lipsum[2]
\end{abstract}
\vspace*{\fill}
\keywords{CO2, carbon dioxide, microcontroller, air quality, CO2 meter}
\addcontentsline{toc}{section}{ABSTRACT}

\newpage
\tableofcontents

\newpage
\listoftables
Tables will be there once they used in document
\addcontentsline{toc}{section}{LIST OF TABLES}

\listoffigures
\addcontentsline{toc}{section}{LIST OF FIGURES}

\newpage
\printglossary[title=LIST OF ABBREVIATIONS, type=\acronymtype]
Use acronyms like that:
its \acrlong{iot}, which is abbreviated \acrshort{iot}. This process 
is similar to that used for the \acrfull{m2m}.
\addcontentsline{toc}{section}{LIST OF ABBREVIATIONS}

\newpage
\setcounter{secnumdepth}{3}
\pagenumbering{arabic}
\section{\centering OVERVIEW}
Begin the first paragraph here.
\par
Begin the second paragraph here.

\subsection{Identification of the need}
Begin the first paragraph here.
\par
Begin the second paragraph here.

\subsection{Definition of the problem}
Begin the first paragraph here.
\par
Begin the second paragraph here.

\subsubsection{Functional requirements}
Begin the first paragraph here.
\par
Begin the second paragraph here.

\subsubsection{Performance requirements}
Begin the first paragraph here.
\par
Begin the second paragraph here.

\subsubsection{Constraints}
Begin the first paragraph here.
\par
Specifically, write about the economic, environmental and social impact of your product. A whole page is dedicated to this in the report guide; please read the guide carefully and write about how these issues apply to your product.

\subsection{Conceptual solutions}
Begin the first paragraph here.
\par
Begin the second paragraph here.

\subsubsection{Literature Review}
Include references to other work and data, for example: Reference \cite{Elawady2020} states that a base station’s power consumption can be divided into four main parts. Bees are flying insects closely related to wasps and ants \cite{muizelaar2007}.
\par
Below is an example of including a figure. You should always give the figure a number and a caption (make sure that this is including in the LIST OF FIGURES) and refer to the figure in the text. \cref{fig:figure1} illustrates the steps electric energy takes from the point of generation to the point of use.
\begin{figure}[H]
    \centering
    \includegraphics[width=\textwidth]{images/1.jpg}
    \caption{Path of electrical energy from the power plant to the home.}
    \label{fig:figure1}
\end{figure}

\subsubsection{Concepts}
Begin the first paragraph here.
\par
Begin the second paragraph here.
\par
Below  is an example of including a table. You should always give the table a number and a caption (make sure that this is including in the LIST OF TABLES) and refer to the table in the text. For example: 
\par
\cref{table:table1} compares different conceptual solutions with respect to the four most important requirements; Concept 2 is chosen for this project due to it’s low cost and reasonable score in all the other categories.\\\\
\begin{table}[H]

\begin{tabular}{ |p{4cm}|p{4cm}|p{3cm}|  }
\hline
\multicolumn{3}{|c|}{Country List} \\
\hline
Country Name or Area Name& ISO ALPHA 2 Code &ISO ALPHA 3 \\
\hline
Afghanistan & AF &AFG \\
Aland Islands & AX   & ALA \\
Albania &AL & ALB \\
Algeria    &DZ & DZA \\
American Samoa & AS & ASM \\
Andorra & AD & AND   \\
Angola & AO & AGO \\
\hline
\end{tabular}

\caption{A table showing some data}
\label{table:table1}

\end{table}

\subsection{Physical architecture}
Begin the first paragraph here.
\par
Begin the second paragraph here.
\begin{figure}[H]
    \centering
    \includegraphics[width=\textwidth]{images/2.jpg}
    \caption{Interface diagram for the system}
    \label{fig:figure2}
\end{figure}
\begin{figure}[H]
    \centering
    \includegraphics[width=0.75\textwidth]{images/3.jpg}
    \caption{Process chart for the system}
    \label{fig:figure3}
\end{figure}

\newpage
\section{\centering WORK PLAN}

\subsection{Work Breakdown Structure}
Begin the first paragraph here.
\par
Begin the second paragraph here.
\begin{figure}[H]
    \centering
    \includegraphics[width=0.75\textwidth]{images/4.jpg}
    \caption{Work Breakdown Structure}
    \label{fig:figure4}
\end{figure}

\subsection{Responsibility Matrix}
Begin the first paragraph here.
\par
Begin the second paragraph here.
\begin{figure}[H]
    \centering
    \includegraphics[width=0.75\textwidth]{images/5.jpg}
    \caption{Responsibility Matrix}
    \label{fig:figure5}
\end{figure}

\subsection{Project Network}
Begin the first paragraph here.
\par
Begin the second paragraph here.
\begin{figure}[H]
    \centering
    \includegraphics[width=0.75\textwidth]{images/6.jpg}
    \caption{Project Network}
    \label{fig:figure6}
\end{figure}

\subsection{Gantt chart}
Since the Gantt chart covers 15 weeks it can be very wide, therefore in this example it is placed on it’s own page and turned sideways. There are various forms for a Gantt chart, choose the form that is suitable for the size and complexity of your project.
\par
Discuss the Gantt chart here (refer to the table).
\par
Begin the first paragraph here.
\par
Begin the second paragraph here.

\subsection{Costs}
Begin the first paragraph here.
\par
Begin the second paragraph here.

\subsection{Risk Assessment}
Begin the first paragraph here.
\par
Begin the second paragraph here.

\newpage
\section{\centering SUB-SYSTEMS}
There will be at least two sub-systems corresponding to two faculty departments and/or work groups.

\subsection{The name of the sub-system 1}
Begin the first paragraph here.
\par
Begin the second paragraph here.

\subsubsection{Requirements}
Begin the first paragraph here.
\par
Begin the second paragraph here.

\subsubsection{Technologies and methods}
Begin the first paragraph here.
\par
Begin the second paragraph here.

\subsubsection{Conceptualization}
Begin the first paragraph here.
\par
Begin the second paragraph here.

\subsubsection{Physical architecture}
Begin the first paragraph here.
\par
Begin the second paragraph here.

\subsubsection{Materialization}
Begin the first paragraph here.
\par
Begin the second paragraph here.

\subsubsection{Evaluation}
Begin the first paragraph here.
\par
Begin the second paragraph here.


\subsection{The name of the sub-system 2}
Begin the first paragraph here.
\par
Begin the second paragraph here.

\subsubsection{Requirements}
Begin the first paragraph here.
\par
Begin the second paragraph here.

\subsubsection{Technologies and methods}
Begin the first paragraph here.
\par
Begin the second paragraph here.

\subsubsection{Conceptualization}
Begin the first paragraph here.
\par
Begin the second paragraph here.

\subsubsection{Physical architecture}
Begin the first paragraph here.
\par
Begin the second paragraph here.

\subsubsection{Materialization}
Begin the first paragraph here.
\par
Begin the second paragraph here.

\subsubsection{Evaluation}
Begin the first paragraph here.
\par
Begin the second paragraph here.

\newpage
\section{\centering INTEGRATION \& EVALUATION}
Begin the first paragraph here.
\par
Begin the second paragraph here.

\subsection{Integration}
Begin the first paragraph here.
\par
Begin the second paragraph here.

\subsection{Evaluation}
Begin the first paragraph here.
\par
Begin the second paragraph here.

\newpage
\section{\centering SUMMARY \& CONCLUSION}
In summary, ...
\par 
Begin the first summary paragraph here.
\par
Begin the second paragraph here.\\\

\par
In conclusion, ...
\par
Begin the first summary paragraph here.
\par
Begin the second paragraph here.

\newpage
\section*{\centering ACKNOWLEDGEMENTS}
We wish to thank our adviser Title Name Surname for ..... . Also acknowledge any other help/support from friends, technicians and other staff etc...
\par
This work was partly/wholly funded by Bahçeşehir University (remove this if are not requesting funding).
\addcontentsline{toc}{section}{ACKNOWLEDGEMENTS}

\newpage
\bibliography{references}

\addcontentsline{toc}{section}{REFERENCES}

\newpage
\section*{APPENDIX A}
Information that does not fit naturally into the main body of the report can be put into an appendix. Typically this would be long sections of software code, product user manuals, large tables of validation results, etc.
\par
An example of providing source code is shown in this appendix. Display the source code in a monospace (fixed-width) font and single-spaced. 
\par
Alternatively give a link to an online code repository.
\par
Example of code listing:
\begin{lstlisting}[language=Python]
import numpy as np
    
def incmatrix(genl1,genl2):
    m = len(genl1)
    n = len(genl2)
    M = None #to become the incidence matrix
    VT = np.zeros((n*m,1), int)  #dummy variable
    
    #compute the bitwise xor matrix
    M1 = bitxormatrix(genl1)
    M2 = np.triu(bitxormatrix(genl2),1) 

    for i in range(m-1):
        for j in range(i+1, m):
            [r,c] = np.where(M2 == M1[i,j])
            for k in range(len(r)):
                VT[(i)*n + r[k]] = 1;
                VT[(i)*n + c[k]] = 1;
                VT[(j)*n + r[k]] = 1;
                VT[(j)*n + c[k]] = 1;
                
                if M is None:
                    M = np.copy(VT)
                else:
                    M = np.concatenate((M, VT), 1)
                
                VT = np.zeros((n*m,1), int)
    
    return M
\end{lstlisting}

\addcontentsline{toc}{section}{APPENDIX A}

\newpage
\section*{APPENDIX B}
Another appendix can go here.
\addcontentsline{toc}{section}{APPENDIX B}

\end{document}
